\section{Axiom 5 (Isometric Covariance)}\label{sctn:axiom-5-(isometric-covariance)}
Now we present the final, and most interesting, axiom in the case of Lorentzian spacetimes.

The interesting aspect of this theorem is understanding what replaces the inhomogeneous Lorentz group connected to the identity, which appears in the Minkowski spacetime analog. One might think to replace it with diffeomorphisms of a Lorentzian spacetime. However, this isn't the correct analog.

Lorentzian spacetimes can admit diffeomorphisms that are not connected to the identity. (For example, in two dimensions \href{https://en.wikipedia.org/wiki/Dehn_twist}{Dehn twists} aren't connected to the identity.) So, if we want to ``mirror'' Axiom 5 (Lorentz Covariance) from the Minkowski case (which only considers group elements connected to the identity), we should limit the diffeomorphisms we employ.

So the ``natural'' choice seems to be the set of diffeomorphisms that are connected to the identity. However, this also isn't the correct analog.

Consider the case of Minkowski spacetime and the analogous axiom Axiom 5 (Lorentz Covariance). There the inhomogeneous Lorentz group connected to the identity is used. Under that group the Minkowski metric is invariant
\begin{align}
    \eta \longrightarrow \eta.
\end{align}
So instead of diffeomorphisms connected to the identity being used, what's being used here are isometries connected to the identity, i.e. diffeomorphism connected to the identity that leave the metric invariant.

So in the current case of a Lorentzian spacetime what we should be using are isometries connected to the identity. Using this, the axiom takes the following form:

(Formalization remark.) An isometry connected to the identity automatically preserves the time orientation, so ``isometries connected to the identity'' is the same group as ``identity-component isometries that preserve the future orientation''. The Lean formalization works with the latter description: it intersects the identity component with the (explicitly orientation-preserving) subgroup, because the inclusion of the identity component into the orientation-preserving isometries relies on a Myers--Steenrod-type rigidity result not yet available in Mathlib. This is an implementation choice only; it does not alter the mathematical content of the axiom below.

\begin{axiom}[Isometric Covariance]
    Let $\mathbf{B}$ be any basis element of the Alexandrov topology on a Lorentzian spacetime $M$, i.e. any set of the form $I^+(p) \cap I^-(q)$.
    
    A member $\varphi$ of the group of isometries of $M$ connected to the identity acts on $\mathfrak{U}(\mathbf{B})$ as follows
    \begin{align}
        \alpha_\varphi : \mathfrak{U}(\mathbf{B}) &\longrightarrow \mathfrak{U}(\varphi(\mathbf{B})) \\
                                   a              &\longmapsto     \alpha_\varphi(a)
    \end{align}
    where $\varphi(\mathbf{B})$ is the image of the basis element $\mathbf{B}$ under the isometry $\varphi$ and $\alpha_\varphi$ is a unital *-isomorphism generated by $\varphi$. The map $\alpha_\varphi$ is such that (1) for the identity isometry $\mathbf{1}$ it satisfies
    \begin{align}
        \alpha_{\mathbf{1}}(a) = a,
    \end{align}
    (2) for all appropriate $a$, $\varphi$, and $\varphi'$ it satisfies
    \begin{align}
        \alpha_{\varphi' \cdot \varphi}(a) = \alpha_{\varphi'}(\alpha_{\varphi}(a)),
    \end{align}
    and (3) for basis elements $\mathbf{B}_\iota \subset \mathbf{B}_\kappa$ and the unital *-monomorphism $i$ of Axiom 2 (Isotony) $\alpha_\varphi$ commutes with $i$. In other words the following diagram

    \begin{center}
        \begin{tikzcd}
            \mathfrak{U}(\mathbf{B}_\kappa) \arrow[r, "\alpha_\varphi"]
            & \mathfrak{U}(\varphi(\mathbf{B}_\kappa)) \\
            \mathfrak{U}(\mathbf{B}_\iota) \arrow[r, "\alpha_\varphi"]  \arrow[u, "i"]
            & \mathfrak{U}(\varphi(\mathbf{B}_\iota))  \arrow[u, "i"]
        \end{tikzcd}
    \end{center}

    commutes.
\end{axiom}
